% GENERATED FILE. Edit canonical lessons, then run npm run generate:books.
% canonical-source-hash: fnv1a64:fa4b9bd7e628d3ef
% canonical-lessons: HI-C391-prayas, HI-C391-parisram, HI-C391-jos, HI-C391-nirantar, HI-C391-dohrana

\chapter{Effort}
\label{ch:hi-effort}
\hlchaptermodality{391}

\begin{chapteropening}
\textbf{By the end of this chapter:} \emph{I can say what effort it will take to succeed.}

Complete the last lesson of chapter 391: I can say what effort it will take to succeed.
\end{chapteropening}

\section[\texorpdfstring{\hi{प्रयास} (prayās)}{प्रयास (prayās)}]{\hi{प्रयास} (prayās) — an effort (formal)}

\label{lesson:HI-C391-prayas}



\begin{quote}
Before the new one: say the Hindi for a step, then the Hindi for to try.
\end{quote}

\subsection*{You'll want to know: \hi{प्रयास}}
\textbf{\hi{प्रयास}} — \emph{prayās} — \textquotedblleft{}an effort (formal)\textquotedblright{}.

\begin{quote}
\hi{अगर} \hi{सब} \hi{मिलकर} \hi{प्रयास} \hi{करेंगे} \hi{तो} \hi{गाँव} \hi{बदलेगा।}
\end{quote}

\begin{quote}
\emph{agar sab milkar prayās kareṅge to gā̃v badlegā.}
\end{quote}

\begin{quote}
If everyone makes an effort together, the village will change.
\end{quote}

Masculine.

\begin{cousinweb}
\textbf{\hi{प्रयास}} (\emph{prayās}) is a tatsama, borrowed from Sanskrit \emph{prayāsa}, \textquotedblleft{}exertion, effort\textquotedblright{} (Monier-Williams).
\end{cousinweb}

\begin{grammarlens}[title={What you've built}]
One more word for saying what happens if.
\end{grammarlens}

\subsection*{Your turn}
\begin{itemize}
\raggedright
  \item \emph{Say it:} \emph{prayās}
  \item \emph{Say it:} \emph{prayās}, once more
  \item \emph{Say it:} \emph{agar sab milkar prayās kareṅge to gā̃v badlegā.} — If everyone makes an effort together, the village will change.
  \item \emph{Recall:} say the Hindi for a step, then the Hindi for to try, then say \emph{prayās} again
\end{itemize}

\begin{tcolorbox}[breakable,colback=teal!4,colframe=teal!35!black,title={Before you move on}]
What does \hi{प्रयास} mean? (\textquotedblleft{}An effort (formal)\textquotedblright{}.) Say it once more.
\end{tcolorbox}
\section[\texorpdfstring{\hi{परिश्रम} (pariśram)}{परिश्रम (pariśram)}]{\hi{परिश्रम} (pariśram) — hard work, toil (formal)}

\label{lesson:HI-C391-parisram}



\begin{quote}
Before the new one: say the Hindi for to try, then the Hindi for an effort.
\end{quote}

\subsection*{You'll want to know: \hi{परिश्रम}}
\textbf{\hi{परिश्रम}} — \emph{pariśram} — \textquotedblleft{}hard work, toil (formal)\textquotedblright{}.

\begin{quote}
\hi{अगर} \hi{तुम} \hi{परिश्रम} \hi{करोगे} \hi{तो} \hi{उसका} \hi{फल} \hi{ज़रूर} \hi{मिलेगा।}
\end{quote}

\begin{quote}
\emph{agar tum pariśram karoge to uskā phal zarūr milegā.}
\end{quote}

\begin{quote}
If you work hard, you'll certainly reap the reward.
\end{quote}

Masculine.

\begin{cousinweb}
\textbf{\hi{परिश्रम}} (\emph{pariśram}) is a tatsama, borrowed from Sanskrit \emph{pariśrama}, \textquotedblleft{}fatigue, exertion, labour\textquotedblright{} (Monier-Williams).
\end{cousinweb}

\begin{grammarlens}[title={What you've built}]
One more word for saying what happens if.
\end{grammarlens}

\subsection*{Your turn}
\begin{itemize}
\raggedright
  \item \emph{Say it:} \emph{pariśram}
  \item \emph{Say it:} \emph{pariśram}, once more
  \item \emph{Say it:} \emph{agar tum pariśram karoge to uskā phal zarūr milegā.} — If you work hard, you'll certainly reap the reward.
  \item \emph{Recall:} say the Hindi for to try, then the Hindi for an effort, then say \emph{pariśram} again
\end{itemize}

\begin{tcolorbox}[breakable,colback=teal!4,colframe=teal!35!black,title={Before you move on}]
What does \hi{परिश्रम} mean? (\textquotedblleft{}Hard work, toil (formal)\textquotedblright{}.) Say it once more.
\end{tcolorbox}
\section[\texorpdfstring{\hi{जोश} (joś)}{जोश (joś)}]{\hi{जोश} (joś) — enthusiasm, fervour}

\label{lesson:HI-C391-jos}



\begin{quote}
Before the new one: say the Hindi for an effort, then the Hindi for toil.
\end{quote}

\subsection*{You'll want to know: \hi{जोश}}
\textbf{\hi{जोश}} — \emph{joś} — \textquotedblleft{}enthusiasm, fervour\textquotedblright{}.

\begin{quote}
\hi{अगर} \hi{टीम} \hi{में} \hi{जोश} \hi{होगा} \hi{तो} \hi{जीत} \hi{पक्की} \hi{है।}
\end{quote}

\begin{quote}
\emph{agar ṭīm meṁ joś hogā to jīt pakkī hai.}
\end{quote}

\begin{quote}
If the team is fired up, victory is certain.
\end{quote}

Masculine.

\begin{cousinweb}
\textbf{\hi{जोश}} (\emph{joś}) is usually traced to Persian \emph{jōsh}, \textquotedblleft{}boiling, ardour\textquotedblright{}.
\end{cousinweb}

\begin{grammarlens}[title={What you've built}]
One more word for saying what happens if.
\end{grammarlens}

\subsection*{Your turn}
\begin{itemize}
\raggedright
  \item \emph{Say it:} \emph{joś}
  \item \emph{Say it:} \emph{joś}, once more
  \item \emph{Say it:} \emph{agar ṭīm meṁ joś hogā to jīt pakkī hai.} — If the team is fired up, victory is certain.
  \item \emph{Recall:} say the Hindi for an effort, then the Hindi for toil, then say \emph{joś} again
\end{itemize}

\begin{tcolorbox}[breakable,colback=teal!4,colframe=teal!35!black,title={Before you move on}]
What does \hi{जोश} mean? (\textquotedblleft{}Enthusiasm, fervour\textquotedblright{}.) Say it once more.
\end{tcolorbox}
\section[\texorpdfstring{\hi{निरंतर} (nirantar)}{निरंतर (nirantar)}]{\hi{निरंतर} (nirantar) — continuous, without a break}

\label{lesson:HI-C391-nirantar}



\begin{quote}
Before the new one: say the Hindi for toil, then the Hindi for enthusiasm.
\end{quote}

\subsection*{You'll want to know: \hi{निरंतर}}
\textbf{\hi{निरंतर}} — \emph{nirantar} — \textquotedblleft{}continuous, without a break\textquotedblright{}.

\begin{quote}
\hi{अगर} \hi{तुम} \hi{निरंतर} \hi{अभ्यास} \hi{करोगे} \hi{तो} \hi{तुम्हारी} \hi{हिंदी} \hi{अच्छी} \hi{हो} \hi{जाएगी।}
\end{quote}

\begin{quote}
\emph{agar tum nirantar abhyās karoge to tumhārī hindī acchī ho jāegī.}
\end{quote}

\begin{quote}
If you practise without a break, your Hindi will become good.
\end{quote}

\begin{cousinweb}
\textbf{\hi{निरंतर}} (\emph{nirantar}) is a tatsama, borrowed from Sanskrit \emph{nirantara}, \textquotedblleft{}without an interval, constant\textquotedblright{} (Monier-Williams).
\end{cousinweb}

\begin{grammarlens}[title={What you've built}]
One more word for saying what happens if.
\end{grammarlens}

\subsection*{Your turn}
\begin{itemize}
\raggedright
  \item \emph{Say it:} \emph{nirantar}
  \item \emph{Say it:} \emph{nirantar}, once more
  \item \emph{Say it:} \emph{agar tum nirantar abhyās karoge to tumhārī hindī acchī ho jāegī.} — If you practise without a break, your Hindi will become good.
  \item \emph{Recall:} say the Hindi for toil, then the Hindi for enthusiasm, then say \emph{nirantar} again
\end{itemize}

\begin{tcolorbox}[breakable,colback=teal!4,colframe=teal!35!black,title={Before you move on}]
What does \hi{निरंतर} mean? (\textquotedblleft{}Continuous, without a break\textquotedblright{}.) Say it once more.
\end{tcolorbox}
\section[\texorpdfstring{\hi{दोहराना} (dohrānā)}{दोहराना (dohrānā)}]{\hi{दोहराना} (dohrānā) — to repeat, to go over again}

\label{lesson:HI-C391-dohrana}



\begin{quote}
Before the new one: say the Hindi for enthusiasm, then the Hindi for without a break.
\end{quote}

\subsection*{You'll want to know: \hi{दोहराना}}
\textbf{\hi{दोहराना}} — \emph{dohrānā} — \textquotedblleft{}to repeat, to go over again\textquotedblright{}.

\begin{quote}
\hi{अगर} \hi{तुम} \hi{यह} \hi{पाठ} \hi{दो} \hi{बार} \hi{दोहराओगे} \hi{तो} \hi{याद} \hi{हो} \hi{जाएगा।}
\end{quote}

\begin{quote}
\emph{agar tum yah pāṭh do bār dohrāoge to yād ho jāegā.}
\end{quote}

\begin{quote}
If you go over this lesson twice, you'll remember it.
\end{quote}

Past, with \hi{ने} (he or she …): \hi{उसने} \hi{सवाल} \hi{दोहराया}, \emph{usne savāl dohrāyā}.

\begin{cousinweb}
\textbf{\hi{दोहराना}} (\emph{dohrānā}) is a verb made on \hi{दोहरा}, \textquotedblleft{}double\textquotedblright{}, which CDIAL traces to a reconstructed \emph{*dudhāra} (CDIAL 6407): to do a thing twice over.
\end{cousinweb}

\begin{grammarlens}[title={What you've built}]
One more word for saying what happens if.
\end{grammarlens}

\subsection*{Your turn}
\begin{itemize}
\raggedright
  \item \emph{Say it:} \emph{dohrānā}
  \item \emph{Say it:} \emph{dohrānā}, once more
  \item \emph{Say it:} \emph{agar tum yah pāṭh do bār dohrāoge to yād ho jāegā.} — If you go over this lesson twice, you'll remember it.
  \item \emph{Recall:} say the Hindi for enthusiasm, then the Hindi for without a break, then say \emph{dohrānā} again
\end{itemize}

\begin{tcolorbox}[breakable,colback=teal!4,colframe=teal!35!black,title={Before you move on}]
What does \hi{दोहराना} mean? (\textquotedblleft{}To repeat, to go over again\textquotedblright{}.) Say it once more.
\end{tcolorbox}
